% 图 4.3.1 几何计算机制 — 基于 CICC0903540 技术文档 4.3.1
\documentclass[tikz,border=8pt]{standalone}
% 顶会/顶刊风格：低饱和度马卡龙 + 高级感
% 适用于 NeurIPS, CVPR, ICLR, IEEE TPAMI
\usepackage{amssymb}
\usepackage{fontspec}
\usepackage{xeCJK}
\setCJKmainfont{SimHei}
\usepackage{tikz}
\usepackage{pgfplots}
\usetikzlibrary{
  arrows.meta, positioning, shapes.geometric, calc,
  backgrounds, fit, decorations.pathreplacing, patterns
}

% 低饱和度马卡龙配色
\definecolor{mBlue}{HTML}{AEC6CF}
\definecolor{mPink}{HTML}{FFD1DC}
\definecolor{mGreen}{HTML}{B1D8B7}
\definecolor{mPurple}{HTML}{B39EB5}
\definecolor{mYellow}{HTML}{FDFD96}
\definecolor{mGray}{HTML}{E5E4E2}
\definecolor{mDark}{HTML}{4A4A4A}

% 全局 TikZ 样式
\tikzset{
  >=stealth,
  every node/.style={align=center, font=\sffamily},
  block/.style={
    rounded corners=3pt, draw=mDark, align=center,
    minimum height=0.8cm, inner sep=5pt},
  arr/.style={-{Stealth[length=4pt]}, semithick, draw=mDark},
  % 信息密度增强：张量维度标注
  ts/.style={font=\tiny\sffamily, color=mDark!80},
  % 数学操作符节点（⊕加 ⊗乘 ‖拼接）
  op/.style={circle, draw=mDark, minimum size=0.4cm, inner sep=0,
    font=\scriptsize\sffamily, fill=mGray!60},
  % 图例
  legendbox/.style={rectangle, draw=mDark!50, rounded corners=2pt,
    fill=mGray!15, inner sep=4pt, font=\tiny\sffamily},
}

\begin{document}
\begin{tikzpicture}[
  x=1.1cm, y=1.0cm,
  node distance=0.6cm and 0.8cm,
  big/.style={
    rectangle, draw=mDark, fill=mBlue!40,
    minimum width=2.2cm, minimum height=0.9cm,
    font=\small\sffamily, align=center, rounded corners=4pt,
    semithick, inner sep=5pt},
  small/.style={
    rectangle, draw=mDark!80, fill=mGreen!50,
    minimum width=1.9cm, minimum height=0.7cm,
    font=\scriptsize\sffamily, align=center, rounded corners=3pt},
  arr/.style={-{Stealth[length=4pt]}, semithick, draw=mDark},
  lbl/.style={font=\scriptsize, color=mGray},
]
  \node[big] (trans) {转换算子\\{\scriptsize\color{mDark!80} 转置、切片}};
  \node[big, right=0.6cm of trans] (comp) {复合算子\\{\scriptsize\color{mDark!80} 卷积、池化}};

  \node[small, below=1.1cm of trans] (raster) {光栅算子\\{\scriptsize\color{mDark!80} 坐标映射}};
  \node[small, below=1.1cm of comp] (atomic) {原子算子\\{\scriptsize\color{mDark!80} 加减乘除}};

  \draw[arr] (trans.south) -- ++(0,-0.4) -| (raster.north);
  \draw[arr] (comp.south) -- ++(0,-0.4) -| (atomic.north);
  \node at ($(trans)!0.5!(comp)+(0,-0.55)$) [lbl] {几何计算分解};

  \node[below=0.35cm of raster, lbl] {优化后 $O(1055)$};
  \node[below=0.35cm of atomic, lbl] {约减少 46\%};

  \begin{scope}[on background layer]
    \node[fit=(trans)(comp), fill=mBlue!25, draw=mDark!50,
      rounded corners=5pt, inner sep=6pt, semithick] {};
    \node[fit=(raster)(atomic), fill=mGreen!35, draw=mDark!50,
      rounded corners=5pt, inner sep=6pt, semithick] {};
  \end{scope}

  % 图例（预留边距）
  \node[legendbox, anchor=north east] at ([xshift=-5pt, yshift=-5pt]current bounding box.south east)
    [align=left, text width=3.5cm] {
    \textbf{图例}\\[1pt]
    \colorbox{mBlue!40}{\rule{0.15cm}{0.1cm}\hspace{0.06cm}} 高层 \quad
    \colorbox{mGreen!50}{\rule{0.15cm}{0.1cm}\hspace{0.06cm}} 底层\\
    region: strides + offsets
  };
\end{tikzpicture}
\end{document}
